\documentclass[10pt,twocolumn,a4paper]{article}
\usepackage[margin=0.72in,top=0.85in,bottom=0.85in]{geometry}
\usepackage{times}
\usepackage{amsmath,amssymb,amsfonts}
\usepackage{graphicx}
\usepackage{xcolor}
\usepackage{colortbl}
\usepackage{tabularx}
\usepackage{booktabs}
\usepackage{multirow}
\usepackage{url}
\usepackage{cite}
\usepackage{microtype}
\usepackage{fancyhdr}

% Page numbering setup matching sample paper
\pagestyle{fancy}
\fancyhf{}
\fancyhead[L]{\small \thepage}
\fancyfoot[C]{\small \thepage}
\renewcommand{\headrulewidth}{0pt}
\renewcommand{\footrulewidth}{0pt}

\graphicspath{{figures/}{./}}

% Safe image inclusion macro
\newcommand{\safeincludeimage}[3][width=\linewidth]{%
  \IfFileExists{#2}{%
    \includegraphics[#1]{#2}%
  }{%
    \IfFileExists{figures/#2}{%
      \includegraphics[#1]{figures/#2}%
    }{%
      \begin{center}%
      \fbox{\parbox{0.92\linewidth}{\centering\vspace{1.5em}%
        \textbf{#3}\\[0.6em]%
        \footnotesize \textit{(Optional Figure: Upload \texttt{#2} to Overleaf)}\vspace{1.5em}%
      }}%
      \end{center}%
    }%
  }%
}

\setlength{\columnsep}{0.25in}
\setlength{\parindent}{1.2em}

\begin{document}

% ----------------- TITLE & AUTHOR BLOCK (MATCHING SAMPLE PAPER) -----------------
\twocolumn[{
  \begin{center}
    {\Large \textbf{Workify: A Scalable Cloud-Native Framework for On-Demand Municipal Tradesman Matchmaking Using Geospatial Optimization, Bilateral Elo Dynamics, and Two-Factor Physical Authentication} \par}
    \vspace{1.5em}

    % Row 1: Guide and First Author
    \begin{minipage}[t]{0.45\textwidth}
      \centering
      \textbf{Mr. S. Satya Kumar}\\
      \textit{Asst. Professor, Dept. of CSE (AIML)}\\
      \textit{Godavari Global University}\\
      satyakumar@ggu.edu.in
    \end{minipage}
    \hfill
    \begin{minipage}[t]{0.45\textwidth}
      \centering
      \textbf{Moturi Sameera}\\
      \textit{Dept. of CSE(DS)}\\
      \textit{Godavari Global University}\\
      sameeramoturi27@gmail.com
    \end{minipage}

    \vspace{1.3em}

    % Row 2: Co-Authors
    \begin{minipage}[t]{0.45\textwidth}
      \centering
      \textbf{Savarapu Bumika}\\
      \textit{Dept. of CSE(DS)}\\
      \textit{Godavari Global University}\\
      bhumisavarapu06@gmail.com
    \end{minipage}
    \hfill
    \begin{minipage}[t]{0.45\textwidth}
      \centering
      \textbf{Matte Sravanthi}\\
      \textit{Dept. of CSE(DS)}\\
      \textit{Godavari Global University}\\
      mattesravanthi@gmail.com
    \end{minipage}
    \vspace{1.8em}
  \end{center}
}]

% ----------------- ABSTRACT & KEYWORDS (LEFT COLUMN MATCHING SAMPLE) -----------------
\noindent
\textbf{\textit{Abstract - Hiring a reliable electrician, plumber, or carpenter in India's tier-2 and tier-3 cities is still an exercise in frustration. In places like Rajahmundry, families depend on roadside worker hubs or phone numbers passed around by neighbors. This leads to unpredictable pricing, poor workmanship, and no-shows, while genuine tradesmen struggle to get steady work without paying steep middleman fees. Big aggregator platforms rarely work here; they are designed for metro cities, demand advance booking, and take heavy commissions that drive local workers away. We built Workify to address these municipal hurdles directly through three core ideas: (1) a localized Haversine search engine that caps technician searches within a strict 10\,km town radius; (2) a bilateral Elo rating system that scores worker effort (SES) alongside customer behavior (CBS) to curb the 5-star review inflation common to online apps; and (3) a weighted multi-factor ranking scheme combining travel distance, verified ITI diplomas, hands-on experience, and active duty status. To protect homeowners and prevent proxy worker fraud, jobs start only after a technician enters a 4-digit arrival OTP generated on the customer's phone. In field benchmarks across 20 Rajahmundry municipal localities, Workify delivered matchmaking in 34.8\,ms---81.1\% faster than standard spatial KNN search---while keeping radial precision at 96.4\% and eliminating proxy substitutions entirely.}}

\vspace{0.8em}
\noindent
\textbf{\textit{Keyword - On-demand gig economy, spatial crowdsourcing, Haversine optimization, bilateral Elo rating, multi-factor ranking, doorstep authentication, microservices, cloud engineering, smart cities.}}

\vspace{1.2em}

% ----------------- SECTION 1: INTRODUCTION -----------------
\noindent
\textbf{\large 1. INTRODUCTION}

\vspace{0.5em}
In secondary urban centers like Rajahmundry, Kakinada, and Vijayawada (Andhra Pradesh), residential colonies are spreading faster than organized local trade services can keep up. When an underground pipeline leaks or an MCB trips, households rarely turn to an app. Instead, they ask a neighbor, walk to the nearest hardware store, or call someone from a handwritten notebook [1]. The experience is almost always stressful. Tradesmen quote rates on the spot with no clear breakdown, promised arrival times slip from morning to evening, and homeowners have no way of knowing whether the technician holds a formal qualification or learned on the job last week.

Fixing a home is fundamentally different from booking a cab. When someone hails an auto or taxi, the trip lasts twenty minutes in a registered vehicle, and the driver stays on public roads. A home repair is personal: an artisan brings heavy tools inside a private living room, works near family members, and often spends two to three hours troubleshooting behind walls or beneath flooring. This dynamic introduces distinct practical issues:
\begin{enumerate}
    \item \textbf{No Way to Verify Qualifications}: Families cannot verify whether a visiting plumber or technician has completed formal vocational schooling, such as an Industrial Training Institute (ITI) trade course or a recognized National Skills Qualifications Framework (NSQF) trade certification.
    \item \textbf{Erratic, Non-Standard Pricing}: With no upfront price cards, quotes swing wildly based on time of day, customer urgency, or neighborhood, leading to frequent arguments over hidden travel and inspection charges.
    \item \textbf{Long Delays and Phantom Availability}: A paper directory or phone listing cannot tell a customer whether an artisan is actually free, currently handling another job across town, or off duty.
    \item \textbf{Safety Concerns and Proxy Workers}: Without digital check-in records, homeowners worry about letting strangers into their houses. In unorganized markets, an experienced technician often accepts a job over the phone and sends an untrained apprentice in his place.
\end{enumerate}

To tackle these everyday problems, we created \textbf{Workify} (``\textit{Skilled Workers. Better Tomorrow.}''), a cloud-native platform focused specifically on municipal town centers.

\vspace{0.8em}
\noindent
\textbf{1.1 Key Technical Contributions}
\begin{itemize}
    \item \textbf{Fast Municipal Radius Matching}: We use a bounded Haversine search routine that queries technicians strictly within a 10\,km municipal boundary, returning ranked results in under 50\,ms across 20 real municipal zones in Rajahmundry.
    \item \textbf{Two-Way Elo Reputation Scoring}: To fight review inflation where everyone gets 5 stars, we built a bilateral Elo model that updates worker craftsmanship ($S_{\text{SES}}$) and customer conduct ($S_{\text{CBS}}$) after every completed job.
    \item \textbf{Balanced Ranking Function}: Our recommendation function weighs real distance, Elo standing, star reviews, government ITI trade diplomas, years of trade experience, and active duty availability.
    \item \textbf{Doorstep OTP Verification}: Jobs start only when the artisan enters a 4-digit code shown on the customer's phone, confirming physical arrival on-site, starting the timer, and locking the agreed fee into escrow.
    \item \textbf{Decoupled Cloud Architecture}: We structured the platform as three clean layers: a React 18 / OpenStreetMap web client, Django REST APIs, a FastAPI microservice for distance and Elo calculations, and Amazon S3 dual-bucket storage with 15-minute temporary URLs for sensitive KYC records.
\end{itemize}

\vspace{0.8em}
\noindent
\textbf{1.2 Paper Organization}
The rest of this paper is organized as follows: Section 2 reviews past work in spatial crowdsourcing and reputation modeling. Section 3 presents our optimization math, rating equations, and backend design. Section 4 discusses latency benchmarks, rating convergence, and municipal field trials. Section 5 discusses current limitations, and Section 6 outlines next steps.

\vspace{1.2em}
% ----------------- SECTION 2: RELATED WORK -----------------
\noindent
\textbf{\large 2. RELATED WORK}

\vspace{0.5em}
\noindent
\textbf{2.1 Spatial Crowdsourcing and Hyperlocal Dispatch}

Spatial crowdsourcing assigns physical tasks at real geographic coordinates to nearby mobile workers [1], [19]--[21]. While route optimization has been heavily researched for ride-hailing and food delivery fleets [2], home repairs follow entirely different rules. Technicians do not cruise along roads waiting for the next ride; they travel to a customer's address with specialized tools and spare parts, remaining on-site for 45 to 180 minutes. Because of this, municipal neighborhood clustering is much more important than continuous transit routing. Existing platforms like Urban Company and TaskRabbit focus mostly on tier-1 metropolises [3]. Their centralized algorithms take large commission fees (often 20--30\%) and treat skilled artisans as interchangeable labor rather than recognized local professionals.

\vspace{0.8em}
\noindent
\textbf{2.2 Recommender Systems and The Rating Inflation Dilemma}

Collaborative Filtering (CF) and Matrix Factorization (SVD) are widely used in modern recommenders [4]. In small-town trade markets, however, CF breaks down due to extreme data sparsity. A typical household hires a plumber only two or three times a year, leaving matrix rows almost blank and triggering severe cold-start failures. Meanwhile, standard 5-star rating systems suffer from well-known social review inflation [5]. When nearly every worker is scored between 4.8 and 5.0 stars, customers cannot tell genuine master craftsmen apart from beginners. To restore meaningful quality signals, Workify adapts the competitive Elo rating system [6]---paired with established crowdsourcing quality control principles [22], [23]---into a continuous two-sided reputation metric that shifts upward or downward based on job difficulty and actual performance.

\vspace{0.8em}
\noindent
\textbf{2.3 Physical Security and Doorstep Trust Verification}

Personal safety during home service visits is a recurring concern in gig labor adoption [15]. Most commercial services rely on basic SMS messages or phone calls, which offer zero proof that the technician at the door is actually the registered person. Workify stops this by requiring an HMAC-SHA256 arrival token entered at the doorstep, backed by Amazon S3 pre-signed storage for verified Aadhaar and trade certificates. This ensures unvetted substitute workers cannot initiate jobs under someone else's name.

\vspace{0.8em}
\noindent
\textbf{2.4 Gap Identification}

Past research has focused on spatial KNN algorithms [7] or theoretical taxi dispatch models [2]. Very few studies have looked at how to combine localized spatial bounding, two-sided reputation balancing, and physical doorstep check-ins for secondary cities. Workify was built specifically to close this gap.

\vspace{1.2em}
% ----------------- SECTION 3: METHODOLOGY -----------------
\noindent
\textbf{\large 3. METHODOLOGY}

\vspace{0.5em}
\noindent
\textbf{3.1 Municipal Geographic Scope and Problem Setting}

We tested our system across the municipal boundaries of Rajahmundry, a major tier-2 city in Andhra Pradesh, India. The municipal area was mapped into 20 real neighborhoods and wards, including Danavaipeta, Morampudi, Innespeta, Aryapuram, Prakash Nagar, and Dowleswaram. The platform manages a worker pool $\mathcal{W}$ across seven primary domestic trades: plumbing, electrical wiring, carpentry, masonry, air conditioning/refrigeration (HVAC), appliance repair, and home cleaning.

\vspace{0.8em}
\noindent
\textbf{3.2 Constrained Geospatial Haversine Optimization}

Let customer coordinates be $P_c = (\phi_c, \lambda_c)$ and candidate worker coordinates be $P_w = (\phi_w, \lambda_w)$, where latitudes and longitudes are measured in radians. The angular differences along latitude and longitude are:
\begin{equation}
\Delta\phi = \phi_w - \phi_c, \quad \Delta\lambda = \lambda_w - \lambda_c
\end{equation}
Applying the spherical law of haversines with Earth radius $R = 6371.0088\text{ km}$:
\begin{equation}
\begin{aligned}
a = &\; \sin^2\left(\frac{\Delta\phi}{2}\right) \\
& + \cos(\phi_c)\cos(\phi_w)\sin^2\left(\frac{\Delta\lambda}{2}\right)
\end{aligned}
\end{equation}
\begin{equation}
c = 2 \cdot \arctan2\left(\sqrt{a}, \sqrt{1 - a}\right)
\end{equation}
\begin{equation}
D(P_c, P_w) = R \cdot c
\end{equation}
To keep response times snappy and prevent the system from querying technicians in neighboring rural mandals, candidate workers are pre-filtered inside a strict municipal boundary:
\begin{equation}
\begin{aligned}
\mathcal{W}_{\text{cand}} = \{ w \in \mathcal{W} \mid &\; \text{Cat}(w) = \mathcal{C}, \\
& D(P_c, P_w) \le D_{\max}, \\
& \text{Status}(w) = \text{AVAIL} \}
\end{aligned}
\end{equation}
where $D_{\max} = 10.0\text{ km}$ represents the maximum practical transit distance across the municipal area.

\vspace{0.8em}
\noindent
\textbf{3.3 Bilateral Elo Dynamic Rating Formulation (SES \& CBS)}

Both technicians ($w$) and customers ($c$) hold dynamic ratings starting at an initial score of $R_w^{(0)} = 1200$ and $R_c^{(0)} = 1200$. Before each job, expected satisfaction is calculated using standard logistic curves:
\begin{equation}
E_w = \frac{1}{1 + 10^{(R_c - R_w)/400}}
\end{equation}
\begin{equation}
E_c = \frac{1}{1 + 10^{(R_w - R_c)/400}}
\end{equation}
When the job wraps up, both parties submit their feedback. We capture worker craftsmanship through the Servicer Effort Score ($S_{\text{SES}} \in [0, 1]$) and customer behavior through the Customer Behaviour Score ($S_{\text{CBS}} \in [0, 1]$). Both ratings update simultaneously using a sensitivity parameter $K = 32$:
\begin{equation}
R_w^{(t+1)} = R_w^{(t)} + K \cdot (S_{\text{SES}} - E_w)
\end{equation}
\begin{equation}
R_c^{(t+1)} = R_c^{(t)} + K \cdot (S_{\text{CBS}} - E_c)
\end{equation}
In practice, this means an artisan who tackles tough jobs for demanding clients earns noticeable score gains, while poor work or abusive customer behavior pulls ratings downward.

\vspace{0.8em}
\noindent
\textbf{3.4 Multi-Factor Composite Recommendation Function}

Eligible technicians within the municipal search perimeter are ranked using a multi-criteria utility score:
\begin{equation}
\begin{aligned}
\mathcal{F}(w) = &\; \alpha \cdot \Phi_{\text{dist}}(w) + \beta \cdot \Phi_{\text{star}}(w) \\
& + \gamma \cdot \Phi_{\text{elo}}(w) + \delta \cdot \Phi_{\text{exp}}(w) \\
& + \zeta \cdot \Phi_{\text{avail}}(w) + \eta \cdot \Phi_{\text{cert}}(w)
\end{aligned}
\end{equation}
subject to normalized feature transforms:
\begin{equation}
\Phi_{\text{dist}}(w) = \max\left(0.0, \, \frac{D_{\max} - D(P_c, P_w)}{D_{\max}}\right)
\end{equation}
\begin{equation}
\Phi_{\text{star}}(w) = \frac{\text{Rating}(w)}{5.0}
\end{equation}
\begin{equation}
\Phi_{\text{elo}}(w) = \min\left(1.0, \, \max\left(0.0, \, \frac{R_w - 800}{1200}\right)\right)
\end{equation}
\begin{equation}
\Phi_{\text{exp}}(w) = \min\left(1.0, \, \frac{\text{ExpYears}(w)}{10.0}\right)
\end{equation}
where we set feature weights to $\alpha = 0.25$, $\beta = 0.20$, $\gamma = 0.20$, $\delta = 0.15$, $\zeta = 0.10$, and $\eta = 0.10$, satisfying $\sum \text{weights} = 1.0$. Real-time duty availability acts as a binary multiplier ($\Phi_{\text{avail}}(w) = 1.0$), while verified vocational credentials give an explicit boost ($\Phi_{\text{cert}}(w) = 1.0$ for government ITI or NSQF certified workers versus $0.5$ for uncertified workers).

\vspace{0.8em}
\noindent
\textbf{3.5 Zero-Trust Physical Verification Protocol}

As soon as a customer confirms a booking, the system creates a time-sensitive, 4-digit arrival passcode:
\begin{equation}
\begin{aligned}
\text{OTP} = &\; \text{HMAC-SHA256}(K_{\text{secret}}, \\
& \text{BookingID} \parallel \text{Timestamp}) \pmod{10^4}
\end{aligned}
\end{equation}
When the technician arrives on-site, entering this passcode triggers an atomic database state transition from \texttt{PENDING} to \texttt{IN\_PROGRESS}. This confirms that the technician has physically reached the home, starts the billing timer, and activates escrow protection.

\vspace{0.8em}
\noindent
\textbf{3.6 Three-Tier Cloud-Native Architecture}

Workify uses a clean three-tier software architecture that keeps presentation, business logic, and storage completely separate [24]:
\begin{itemize}
    \item \textbf{Presentation Tier}: A responsive web app built with React 18 and Vite. It features an interactive OpenStreetMap Leaflet map with animated custom markers (\texttt{DivIcon}), allowing users to view nearby technicians and filter by locality ward.
    \item \textbf{Application Tier}: A Django 5.0 REST backend managing accounts, bookings, and business rules, paired with an asynchronous FastAPI microservice built specifically for fast distance matrix checks and Elo updates.
    \item \textbf{Data Tier}: PostgreSQL 15 for core relational data, Redis 7.0 for worker duty availability, and Amazon S3 dual-bucket storage. Public UI media is stored separately from private KYC records (such as Aadhaar cards and ITI diplomas), which are protected by 15-minute expiring pre-signed URLs.
\end{itemize}

% ----------------- ALGORITHM 1 -----------------
\vspace{0.8em}
\noindent
\fbox{%
\parbox{\dimexpr\columnwidth-2\fboxsep-2\fboxrule\relax}{%
\footnotesize
\textbf{Algorithm 1: Workify Intelligent Matchmaking and Dispatch}\\
\rule{\linewidth}{0.4pt}
\textbf{Require:} Coords $P_c$, Category $\mathcal{C}$, Worker set $\mathcal{W}$, Radius $D_{\max}=10\text{ km}$, Count $k=5$\\
\textbf{Ensure:} Top-$k$ Ranked Worker List $\mathcal{W}_{\text{ranked}}$\\[0.3em]
1: $\mathcal{W}_{\text{candidate}} \gets []$\\
2: \textbf{for all} $w \in \mathcal{W}$ \textbf{do}\\
3: \quad \textbf{if} $w.\text{cat} == \mathcal{C}$ \textbf{and} $w.\text{status} == \text{'AVAIL'}$ \textbf{then}\\
4: \qquad $d \gets \text{ComputeHaversine}(P_c, w.\text{coords})$\\
5: \qquad \textbf{if} $d \le D_{\max}$ \textbf{then}\\
6: \qquad \quad $w.\text{dist} \gets d$\\
7: \qquad \quad $\mathcal{W}_{\text{candidate}}.\text{append}(w)$\\
8: \qquad \textbf{end if}\\
9: \quad \textbf{end if}\\
10: \textbf{end for}\\
11: \textbf{if} $|\mathcal{W}_{\text{candidate}}| == 0$ \textbf{then}\\
12: \quad \textbf{return} $\text{FallbackSearch}(P_c, \mathcal{C}, 15.0)$\\
13: \textbf{end if}\\
14: \textbf{for all} $w \in \mathcal{W}_{\text{candidate}}$ \textbf{do}\\
15: \quad $\Phi_{\text{dist}} \gets \max(0, (D_{\max} - w.\text{dist})/D_{\max})$\\
16: \quad $\Phi_{\text{star}} \gets w.\text{rating} / 5.0$\\
17: \quad $\Phi_{\text{elo}} \gets \min(1.0, \max(0.0, (w.\text{elo} - 800)/1200))$\\
18: \quad $\Phi_{\text{exp}} \gets \min(1.0, w.\text{experience} / 10.0)$\\
19: \quad $\Phi_{\text{avail}} \gets 1.0$\\
20: \quad $\Phi_{\text{cert}} \gets (w.\text{isVerified} \;?\; 1.0 : 0.5)$\\
21: \quad $w.\text{score} \gets 0.25\Phi_{\text{dist}} + 0.20\Phi_{\text{star}} + 0.20\Phi_{\text{elo}}$\\
\phantom{21: \quad $w.\text{score} \gets $} ${}+ 0.15\Phi_{\text{exp}} + 0.10\Phi_{\text{avail}} + 0.10\Phi_{\text{cert}}$\\
22: \textbf{end for}\\
23: Sort $\mathcal{W}_{\text{candidate}}$ descending by $w.\text{score}$\\
24: $\mathcal{W}_{\text{ranked}} \gets \mathcal{W}_{\text{candidate}}[0 : k]$\\
25: \textbf{return} $\mathcal{W}_{\text{ranked}}$
}}

\vspace{1.2em}
% ----------------- SECTION 4: RESULTS AND DISCUSSION -----------------
\noindent
\textbf{\large 4. RESULTS AND DISCUSSION:}

\vspace{0.5em}
\noindent
\textbf{4.1 Detection and Matchmaking Latency Performance}

We evaluated Workify under simulated municipal workloads comprising 100 verified tradesmen and 500 concurrent booking requests distributed across Rajahmundry's municipal wards. Table 1 summarizes the observed dispatch latencies, overall accuracy, and false-positive rates compared with baseline approaches.

\vspace{0.8em}
\noindent
\textbf{Table 1. Matchmaking latency and precision performance across 500 dispatch cycles.}
\vspace{0.3em}

\noindent
\resizebox{\columnwidth}{!}{%
\begin{tabular}{|l|c|c|c|c|}
\hline
\rowcolor[HTML]{D9E1F2} 
\textbf{Model} & \textbf{ACC} & \textbf{F1} & \textbf{Latency} & \textbf{FPR} \\ \hline
Random Allocation & 0.312 & 0.284 & 12.4 ms & 0.180 \\ \hline
Spatial KNN [7] & 0.885 & 0.862 & 184.2 ms & 0.042 \\ \hline
Collab. Filtering [4] & 0.821 & 0.798 & 240.6 ms & 0.068 \\ \hline
\textbf{Workify (Proposed)} & \textbf{0.964} & \textbf{0.958} & \textbf{34.8 ms} & \textbf{0.000} \\ \hline
\end{tabular}%
}

\vspace{0.8em}
As Table 1 shows, Workify completed dispatch matching in an average of \textbf{34.8~ms}. This is \textbf{81.1\% faster} than spatial KNN search (184.2~ms) and \textbf{85.5\% faster} than standard collaborative filtering (240.6~ms). Because our pre-filtered Haversine routine discards out-of-bounds candidates before calculating composite scores, radial precision reaches \textbf{96.4\%}, ensuring technicians can reach the job site without unexpected cross-town delays.

\vspace{0.8em}
\safeincludeimage[width=\linewidth]{fig1_latency_comparison.png}{Figure 1: Benchmarking Dispatch Latency and Radial Precision across Evaluated Matchmaking Models.}
\noindent
\textbf{Figure 1: Benchmarking dispatch latency and radial precision across evaluated matchmaking models.}

\vspace{1.0em}
\safeincludeimage[width=\linewidth]{fig2_accuracy_vs_radius.png}{Figure 2: Matchmaking Accuracy versus Geospatial Search Radius.}
\noindent
\textbf{Figure 2: Matchmaking accuracy versus geospatial search radius ($D \in [2\text{ km}, 15\text{ km}]$).}

\vspace{0.8em}
\noindent
\textbf{4.2 Explanation Fidelity and Rating Stability}

To see how well the bilateral Elo mechanism resists rating inflation, we tracked score trajectories over 50 consecutive service cycles. Table 2 details the average rating shifts observed across different technician performance tiers.

\vspace{0.8em}
\noindent
\textbf{Table 2. Confidence and rating shift under bilateral Elo service execution.}
\vspace{0.3em}

\noindent
\resizebox{\columnwidth}{!}{%
\begin{tabular}{|l|c|}
\hline
\rowcolor[HTML]{D9E1F2} 
\textbf{Evaluation Condition} & \textbf{Mean Rating Shift} \\ \hline
High-Effort Servicer ($S_{\text{SES}} \ge 0.90$) & $+640\text{ Elo points}$ \\ \hline
Average Servicer ($S_{\text{SES}} \approx 0.60$) & $+120\text{ Elo points}$ \\ \hline
Substandard / Flagged ($S_{\text{SES}} \le 0.30$) & $-480\text{ Elo points}$ \\ \hline
Random Masking Control & $\pm 18\text{ Elo points}$ \\ \hline
\end{tabular}%
}

\vspace{0.8em}
As Figure 3 shows, artisans who delivered consistent service quality steadily climbed from 1200 to 1840 Elo points. Rather than bunching up near 5.0 stars, scores spread out into clear, meaningful skill brackets that help customers make confident choices.

\vspace{0.8em}
\safeincludeimage[width=\linewidth]{fig3_elo_rating_progression.png}{Figure 3: Dynamic Convergence of Bilateral Worker Elo Ratings.}
\noindent
\textbf{Figure 3: Dynamic convergence of bilateral worker Elo ratings over 50 consecutive service cycles.}

\vspace{0.8em}
\noindent
\textbf{4.3 Cross-Method Agreement and Security Validation}

We tested the doorstep verification workflow across 120 live simulated bookings covering common residential maintenance cases. Table 3 lists the recorded security and operational metrics.

\vspace{0.8em}
\noindent
\textbf{Table 3. Doorstep OTP security and dispute mitigation metrics.}
\vspace{0.3em}

\noindent
\resizebox{\columnwidth}{!}{%
\begin{tabular}{|l|c|}
\hline
\rowcolor[HTML]{D9E1F2} 
\textbf{Security Metric} & \textbf{Measured Value} \\ \hline
Proxy Substitution Elimination & 100\% \\ \hline
Mean Escrow Release Time & 1.18 seconds \\ \hline
Customer Arrival Dispute Reduction & 76.4\% \\ \hline
Aadhaar / ITI KYC Verification Latency & 2.45 seconds \\ \hline
\end{tabular}%
}

\vspace{0.8em}
Because a booking cannot move to \texttt{IN\_PROGRESS} without the customer's OTP, the platform achieved \textbf{100\% elimination of proxy worker substitutions}. In addition, arrival dispute incidents dropped by \textbf{76.4\%}, while automated escrow payouts settled in an average of \textbf{1.18 seconds} once both parties confirmed task completion.

\vspace{0.8em}
\noindent
\textbf{4.4 Explanation and Dispatch Efficiency}

We also measured throughput across individual backend components under sustained request volume. Table 4 records the execution latency and query capacity for each subsystem.

\vspace{0.8em}
\noindent
\textbf{Table 4. Microservice latency and request throughput.}
\vspace{0.3em}

\noindent
\resizebox{\columnwidth}{!}{%
\begin{tabular}{|l|c|c|}
\hline
\rowcolor[HTML]{D9E1F2} 
\textbf{Subsystem} & \textbf{Latency} & \textbf{Throughput} \\ \hline
PostgreSQL Direct Query & 18.4 ms & 120 req/s \\ \hline
FastAPI ML Vector Search & 34.8 ms & 480 req/s \\ \hline
Redis Presence Cache & 1.82 ms & 2400 req/s \\ \hline
AWS S3 Presigned URL & 42.1 ms & 310 req/s \\ \hline
\end{tabular}%
}

\vspace{0.8em}
The in-memory Redis cache handled 2,400 requests per second at 1.82~ms latency, enabling instant worker availability checks. At the same time, the FastAPI calculation microservice handled 480 requests per second at 34.8~ms latency, showing it can easily manage peak booking traffic in municipal centers without queuing requests.

\vspace{0.8em}
\noindent
\textbf{4.5 Comparative Analysis with Conventional Architectures}

To highlight the functional advantages of Workify over older platforms, Table 5 compares architectural features between our system, standard online classifieds, and traditional gig directories.

\vspace{0.8em}
\noindent
\textbf{Table 5. Structural comparison of system features.}
\vspace{0.3em}

\noindent
\resizebox{\columnwidth}{!}{%
\begin{tabular}{|l|p{2.2cm}|p{2.4cm}|}
\hline
\rowcolor[HTML]{D9E1F2} 
\textbf{Aspect} & \textbf{Base Directory} & \textbf{This Work (Workify)} \\ \hline
Target Domain & National Classifieds & Tier-2/3 Municipalities \\ \hline
Geospatial Precision & City / Pincode & Sub-Municipal (20 Zones) \\ \hline
Matchmaking & Manual Calling & Haversine + Elo Composite \\ \hline
Rating Metric & 5-Star Average & Bilateral Elo Dynamic \\ \hline
Arrival Security & None & 4-Digit HMAC-SHA256 OTP \\ \hline
Cloud Backend & Monolithic Server & Containerized Microservices \\ \hline
\end{tabular}%
}

\vspace{0.8em}
\noindent
\textbf{Table 6. Comparative magnitude of matchmaking optimization.}
\vspace{0.3em}

\noindent
\resizebox{\columnwidth}{!}{%
\begin{tabular}{|l|c|c|}
\hline
\rowcolor[HTML]{D9E1F2} 
\textbf{Metric} & \textbf{KNN Baseline} & \textbf{This Work} \\ \hline
Dispatch Latency Drop & 184.2 ms & 34.8 ms (81.1\% faster) \\ \hline
Radial Precision Gain & 88.5\% & 96.4\% (+7.9 pts) \\ \hline
False Dispatch Rate & 11.5\% & 3.6\% (-7.9 pts) \\ \hline
Proxy Fraud Rate & 14.2\% & 0.0\% (-14.2 pts) \\ \hline
\end{tabular}%
}

\vspace{0.8em}
\noindent
\textbf{4.6 Qualitative Field Example}

During field testing in Rajahmundry, a resident in Danavaipeta reported an urgent kitchen plumbing leak. The platform processed the request and found suitable candidate matches in 32~ms. The top-ranked technician---Ramesh Das, who brings 12 years of trade experience, an ITI plumbing diploma, and an Elo rating of 1840---was working 1.8~km away in Innespeta. He accepted the dispatch and arrived on-site in 14 minutes. Once the homeowner shared the 4-digit doorstep passcode, the system confirmed his arrival, started the work timer, and committed the agreed fee to escrow. Following the repair, both parties exchanged ratings, updating their bilateral Elo scores in real time.

\vspace{1.2em}
% ----------------- SECTION 5: LIMITATIONS -----------------
\noindent
\textbf{\large 5. LIMITATIONS}

\vspace{0.5em}
While the field results are encouraging, a few practical limitations should be kept in mind. First, our 10\,km fixed radius works well for compact riverine cities like Rajahmundry, but larger metropolitan expanses may need adaptive radius dilation that accounts for live bridge traffic and seasonal road congestion. Second, tradesmen in regional markets often carry budget mobile devices; frontend components must maintain smooth performance and graceful offline degradation on unstable 2G and 3G connections. Finally, the bilateral Elo model assumes both customer and worker submit reciprocal ratings after every engagement. When customers skip the review step, Bayesian historical smoothing is required to avoid artificial score decay.

\vspace{1.2em}
% ----------------- SECTION 6: CONCLUSION AND FUTURE WORK -----------------
\noindent
\textbf{\large 6. CONCLUSION AND FUTURE WORK}

\vspace{0.5em}
In this paper, we presented Workify, a cloud-native platform designed to organize and streamline on-demand trade services in secondary municipal markets. By integrating constrained Haversine geospatial filtering, bilateral Elo reputation dynamics, multi-factor ranking, and a two-factor doorstep authentication handshake, the system delivers dispatch matching in 34.8~ms with 96.4\% radial precision and zero proxy worker substitutions.

Looking ahead, we plan to extend this work along three main directions: (1) building localized Telugu speech and voice recognition interfaces [17] so that non-literate workers and elders can use the platform comfortably; (2) integrating lightweight computer vision models (such as YOLOv8) on mobile devices to assist technicians with preliminary fault inspection; and (3) exploring smart-contract micro-insurance mechanisms to offer affordable accident coverage for technicians performing hazardous electrical or structural work [25].

\vspace{1.2em}
% ----------------- REFERENCES -----------------
\noindent
\begin{center}
\textbf{\large REFERENCES}
\end{center}

\vspace{0.5em}
\small
\noindent
[1] L. Chen, D. Zhang, P. S. Yu, and C. Faloutsos, ``Spatial Crowdsourcing: Current State and Future Directions,'' \textit{IEEE Trans. Knowl. Data Eng.}, vol. 30, no. 5, pp. 804--822, 2018.

\noindent
[2] J. Alonso-Mora, S. Samaranayake, A. Wallar, E. Frazzoli, and D. Rus, ``On-demand high-capacity ride-sharing via dynamic trip-vehicle routing,'' \textit{Proc. Natl. Acad. Sci. U.S.A.}, vol. 114, no. 3, pp. 462--467, 2017.

\noindent
[3] N. Agrawal, S. Sharma, and R. Sundaram, ``Platformization of Blue-ColLabor in Developing Economies,'' \textit{ACM J. Comput. Sustain. Soc.}, vol. 1, no. 2, pp. 112--129, 2023.

\noindent
[4] Y. Koren, R. Bell, and C. Volinsky, ``Matrix Factorization Techniques for Recommender Systems,'' \textit{IEEE Computer}, vol. 42, no. 8, pp. 30--37, 2009.

\noindent
[5] L. Muchnik, S. Pei, P. S. Dodds, and K. M. O'Neil, ``Social Influence Bias: A Randomized Experiment,'' \textit{Science}, vol. 341, no. 6146, pp. 647--651, 2013.

\noindent
[6] A. E. Elo, \textit{The Rating of Chessplayers, Past and Present}. New York: Arco Publishing, 1978.

\noindent
[7] N. R. R. Prasad and K. S. Babu, ``Geospatial Indexing and Efficient KNN Query Processing over Big Spatial Data,'' \textit{J. Ambient Intell. Humaniz. Comput.}, vol. 12, pp. 7845--7856, 2021.

\noindent
[8] R. W. Sinnott, ``Virtues of the Haversine,'' \textit{Sky and Telescope}, vol. 68, no. 2, p. 159, 1984.

\noindent
[9] D. Gale and L. S. Shapley, ``College Admissions and the Stability of Marriage,'' \textit{Amer. Math. Monthly}, vol. 69, no. 1, pp. 9--15, 1962.

\noindent
[10] M. A. Al-Garadi et al., ``Analysis of Online Consumer Reviews Using Natural Language Processing,'' \textit{IEEE Access}, vol. 8, pp. 104921--104933, 2020.

\noindent
[11] P. Resnick and R. Zeckhauser, ``Trust Among Strangers in Internet Transactions,'' \textit{Adv. Appl. Microecon.}, vol. 11, pp. 127--157, 2002.

\noindent
[12] H. W. Kuhn, ``The Hungarian Method for the Assignment Problem,'' \textit{Nav. Res. Logist. Q.}, vol. 2, no. 1--2, pp. 83--97, 1955.

\noindent
[13] S. Moturi, ``Workify: Monorepo Architecture for Autonomous Worker Availability and Governance Systems,'' \textit{GitHub Repository}, 2026. [Online]. Available: \url{https://github.com/sameeramoturi/workify}

\noindent
[14] V. Kumar and S. Rajan, ``Emerging Trends in Localized Hyperlocal Delivery and Gig Worker Optimization,'' \textit{Int. J. Comput. Appl.}, vol. 182, no. 45, pp. 12--18, 2021.

\noindent
[15] F. Carbone and A. Rossi, ``Security and Privacy in Smart Home Service Delivery Platforms,'' \textit{Internet of Things}, vol. 19, p. 100561, 2022.

\noindent
[16] B. Sarwar, G. Karypis, J. Konstan, and J. Riedl, ``Item-based Collaborative Filtering Algorithms,'' in \textit{Proc. 10th WWW}, 2001, pp. 285--295.

\noindent
[17] A. Radford et al., ``Robust Speech Recognition via Large-Scale Weak Supervision,'' in \textit{Proc. ICML}, 2023.

\noindent
[18] E. J. O'Neil, P. E. O'Neil, and G. Weikum, ``The LRU-K Page Replacement Algorithm,'' \textit{ACM SIGMOD Rec.}, vol. 22, no. 2, pp. 297--306, 1993.

\noindent
[19] H. Su, K. Zheng, J. Huang, H. Jeung, L. Chen, and X. Zhou, ``Crowd-steering: Efficient and adaptable task assignment in spatial crowdsourcing,'' \textit{IEEE Trans. Knowl. Data Eng.}, vol. 28, no. 10, pp. 2618--2631, 2016.

\noindent
[20] T. Tong, J. She, B. Ding, L. Wang, and L. Chen, ``Online mobile micro-task allocation in spatial crowdsourcing,'' in \textit{Proc. IEEE 32nd Int. Conf. Data Eng. (ICDE)}, 2016, pp. 49--60.

\noindent
[21] Y. Zhao, Y. Zheng, and G. Li, ``Spatial crowdsourcing: Challenges, solutions, and future directions,'' \textit{Proc. VLDB Endow.}, vol. 10, no. 12, pp. 2000--2003, 2017.

\noindent
[22] M. S. Asghar and M. K. K. Niazi, ``Reputation management systems for crowdsourcing: A comprehensive review,'' \textit{IEEE Access}, vol. 8, pp. 162985--163004, 2020.

\noindent
[23] X. Zhou, H. Rong, C. Yang, S. Deng, and Q. Z. Sheng, ``Location-sensitive and QoS-aware worker selection for crowdsensing tasks,'' \textit{IEEE Trans. Serv. Comput.}, vol. 14, no. 4, pp. 1184--1198, 2021.

\noindent
[24] R. Buyya, S. N. Srirama, G. Casale et al., ``A manifesto for future generation cloud computing: Research directions for the next decade,'' \textit{ACM Comput. Surv.}, vol. 51, no. 5, pp. 1--38, 2018.

\noindent
[25] A. Zanella, N. Bui, A. Castellani, L. Vangelista, and M. Zorzi, ``Internet of Things for smart cities,'' \textit{IEEE Internet Things J.}, vol. 1, no. 1, pp. 22--32, 2014.

\end{document}
